\section{Zaključak}

\paragraph{Zadatak 1}
Problem je konveksan: Hesseova matrica ima svojstvene vrijednosti $1.955969$ i $4.044031$, dakle $Q \succ 0$, a ograničenja su afina. Od četiri kombinacije aktivnosti ograničenja KKT uvjete zadovoljava samo ona s aktivnim $g_1$, što daje
\begin{equation*}
x^* = \left( \tfrac{190}{33},\ \tfrac{140}{33} \right) = (5.757576,\ 4.242424), \qquad
\mu^* = \left( \tfrac{404}{33},\ 0 \right), \qquad
p^* = \tfrac{2000}{33} = 60.606061 .
\end{equation*}
Rješavač \texttt{quadprog} vraća iste primarne i dualne varijable. Dualni problem je opet kvadratni program u $\mu$; njegovo rješenje $\mu^*$ podudara se s množiteljima iz KKT uvjeta, a dualni procjep iznosi $2.48 \cdot 10^{-10}$, pa jaka dualnost vrijedi, kako i predviđa zadovoljen Slaterov uvjet. Optimalna vrijednost osjetljiva je samo na perturbaciju prvog ograničenja, $\partial p^* / \partial v_1 = -12.242424$, dok je $\partial p^* / \partial v_2 = 0$ jer je $g_2$ neaktivno.

\paragraph{Zadatak 2}
Uz nesigurnost koeficijenata od $\pm 15\%$ svako se robusno ograničenje svodi na četiri nejednakosti u vrhovima pravokutnog skupa nesigurnosti, pa robusni problem postaje kvadratni program s osam linearnih nejednakosti. Njegovo rješenje je
\begin{equation*}
x_{\mathrm{rob}}^* = \left( \tfrac{13}{2},\ \tfrac{179}{34} \right) = (6.500000,\ 5.264706), \qquad
f(x_{\mathrm{rob}}^*) = 84.623962 ,
\end{equation*}
uz dva aktivna vrha i strogo pozitivne pripadajuće množitelje $\mu = (17.607439,\ 1.344118)$. U uzorku od 100 slučajnih realizacija ($\texttt{rng(42)}$) nominalno rješenje krši ograničenja u 54 slučaja, a robusno ni u jednom; dopustivost na cijelom skupu nesigurnosti jamči zadovoljenje svih osam vršnih nejednakosti. Za tu se dopustivost plaća porast funkcije cilja od $24.017901$, odnosno $39.63\%$.

\paragraph{Zadatak 3}
Nakon diskretizacije konačnim razlikama i uvođenja pomoćnih varijabli $z_i \geq |y_i - h_i|$ problem pripada klasi \textbf{linearnog programiranja}, s $2(N+1)$ varijabli i $6N$ nejednakosti. Rezultati za četiri zadana slučaja jesu:
\begin{equation*}
\begin{array}{lcc}
\text{slučaj A}, N = 20:  & J^* = 1.916180\ \mathrm{m}^2 \qquad & \text{slučaj B}, N = 20:  \ J^* = 2.453727\ \mathrm{m}^2 \\
\text{slučaj A}, N = 100: & J^* = 1.965940\ \mathrm{m}^2 \qquad & \text{slučaj B}, N = 100: \ J^* = 2.506047\ \mathrm{m}^2
\end{array}
\end{equation*}
Gušća mreža povećava vrijednost kriterija za oko $2\%$ jer pravokutna kvadratura na rijetkoj mreži podcjenjuje integral. Pooštravanje ograničenja brzine promjene nagiba s $b_2 = 0.5$ na $b_2 = 0.2\ \mathrm{m}^{-1}$ zahtijeva glađu niveletu i povećava volumen zemljanih radova za $27.47\%$; u oba slučaja aktivna su i ograničenje nagiba i ograničenje brzine promjene nagiba.

\paragraph{Zadatak 4}
Obje matrice modova su Hurwitzove, sa svojstvenim vrijednostima $-1 \pm j\pi/2$. Unatoč tome \textbf{sklopni sustav nije stabilan}: matrica prijelaza preko jednog perioda prespajanja iznosi $\Phi_{\mathrm{cikl}} = e^{-2}\,\mathrm{diag}(-9,\ -1/9)$, pa je $\rho(\Phi_{\mathrm{cikl}}) = 9e^{-2} = 1.218018 > 1$ i norma stanja raste, što potvrđuje i simulacija. Zahtjev da postoji zajednička kvadratna Ljapunovljeva funkcija preveden je kongruentnom transformacijom i supstitucijom $\mathcal{Q} = P^{-1}$, $Y = KP^{-1}$ u sustav linearnih matričnih nejednakosti. Njegovim rješavanjem dobiven je regulator
\begin{equation*}
K = \begin{bmatrix} -3.424629 & -0.104650 \end{bmatrix},
\end{equation*}
uz $P \succ 0$ za koji je $\lambda_{\max}\big((A_i + BK)^{\mathsf{T}} P + P(A_i + BK)\big)$ jednak $-0.910478$ odnosno $-0.329090$ za dva moda. Zatvoreni krug je time eksponencijalno stabilan uz zajamčenu stopu $\alpha = 0.288374\ \mathrm{s}^{-1}$, a simulacija pokazuje pad norme stanja na $1.52 \cdot 10^{-3}$ unutar 5 sekundi.
